\section{Solution}
We decomposed the problem of estimating the ski jumper’s trajectory into a sequence of well-defined sub-tasks, so that each component could be designed, tested, and improved independently.

\medskip

\textbf{Athlete detection (YOLO) and ROI extraction.}  
For each frame we detect the athlete using YOLO. The resulting bounding box is used to define a region of interest (ROI) that focuses computation on the jumper area and reduces sensitivity to irrelevant image content. This step also provides a consistent input for the downstream segmentation and tracking modules.

\medskip

\textbf{Pixel-level segmentation (SAM3) and center-of-mass measurement.}  
Given the ROI, we apply SAM3 to obtain a pixel-accurate segmentation mask of the jumper. From this binary mask we compute the athlete’s \textbf{center of mass (CoM)} by averaging the pixel coordinates belonging to the foreground region. The CoM is used as the primary measurement to represent the jumper position over time and produces an initial (raw) 2D trajectory in image coordinates.

\medskip

\textbf{Temporal stabilization (Kalman filter).}  
Raw CoM measurements can exhibit jitter and occasional outliers due to segmentation noise, partial occlusions, and motion blur. To enforce temporal coherence, we apply a Kalman filter to the CoM sequence (constant-velocity motion model). This provides a smoother and more reliable stream of observations while preserving the overall dynamics of the motion, and it reduces the impact of isolated bad measurements.

\medskip

\textbf{Shot segmentation (HSV cut detection).}  
Broadcast ski-jumping videos are multi-camera and contain abrupt viewpoint changes. Since geometric alignment is only meaningful within a consistent viewpoint, we explicitly detect camera switches. We use an HSV-based method that flags frames where the global color distribution changes sharply, which is a strong indicator of a cut. The video is then split into homogeneous segments (\textbf{shots}) that are processed independently.

\medskip

\textbf{Per-shot motion compensation (feature-based alignment).}  
For each shot, we choose as reference the \textbf{first frame in which the jumper is visible} and express all subsequent measurements in this reference coordinate system. To compensate for camera motion, we estimate the transformation between \textbf{consecutive frames} using a feature-based pipeline:

\begin{itemize}
    \item \textbf{Grayscale conversion and background masking.}  
    Frames are converted to grayscale for feature extraction. To avoid bias from the moving athlete, we restrict feature detection to the \textbf{static background} using a valid-pixel mask (derived from the segmentation), so that the estimated motion reflects camera motion rather than object motion.

    \item \textbf{SIFT keypoints and descriptors.}  
    We extract SIFT keypoints and descriptors on the masked images. SIFT is well-suited to this scenario because it is relatively robust to scale changes and moderate viewpoint changes.

    \item \textbf{Descriptor matching with L2 norm + Lowe ratio test.}  
    SIFT descriptors are real-valued vectors, therefore we match them using the \textbf{Euclidean distance (L2 norm)} in descriptor space. However, raw nearest-neighbor matching is often unreliable in scenes with repeated textures (e.g., snow patterns, stands).  
    To reject ambiguous correspondences, we perform \textbf{k matching with $k=2$} and apply the \textbf{Lowe ratio test}: a match is accepted only if the best match is significantly closer than the second-best one (in our implementation, threshold $0.75$). This filters out many unstable matches and improves robustness before geometric estimation.
\end{itemize}

\medskip

\textbf{Robust transform estimation with RANSAC (affine partial 2D).}  
Even after the ratio test, some outliers remain (mismatches, dynamic elements, blur). We therefore estimate the inter-frame transformation using \textbf{RANSAC} inside \texttt{estimateAffinePartial2D}. RANSAC repeatedly:
(1) samples a minimal subset of correspondences,  
(2) fits a candidate model,  
(3) classifies inliers using a reprojection error threshold (3 pixels in our setup), and  
(4) selects the model with the largest consensus set.  
This makes the estimation robust to outliers and prevents a small number of wrong matches from corrupting the motion model. We additionally validate the solution by requiring a minimum number of inliers and a minimum inlier ratio; otherwise we fall back to the identity transform for that step.

\medskip

We adopt an \textbf{affine partial 2D} model (translation + rotation + isotropic scale, no shear). Compared to a full homography, this model is more stable in our sequences, where motion blur and limited texture can cause projective estimation to overfit and drift.

\medskip

\textbf{Cumulative composition to the shot reference.}  
Let $W_{k\rightarrow k-1}$ be the affine partial transform mapping coordinates from frame $k$ to frame $k-1$. We build the mapping from each frame to the reference frame $0$ via cumulative composition:
\[
H_{0\rightarrow 0} = I, \qquad
H_{k\rightarrow 0} = W_{k\rightarrow k-1}\, H_{k-1\rightarrow 0}.
\]
Given the filtered CoM measurement $p_k = (x_k, y_k)$, we transform it into the reference coordinate system using homogeneous coordinates:
\[
\tilde{p}_k \sim H_{k\rightarrow 0}\,[x_k \; y_k \; 1]^T.
\]
This produces a trajectory that is consistent within each shot and compensated for camera motion.

\medskip

\textbf{Final output.}  
The final result is a set of per-shot trajectories expressed in a common reference frame for each shot.

\section{From 2D to 3D}
Once the 2D detections have been stabilized, all trajectory points are expressed in a single
\emph{reference frame} (the chosen keyframe) by composing the inter-frame warps.
In this way, the motion is no longer described in a moving camera coordinate system, but in a
common image coordinate system where geometric reasoning becomes consistent.

\medskip

Our goal is then to remove the remaining \emph{perspective distortion} by reasoning directly on the
reference frame. Since we assume that the jumper trajectory lies on a single physical plane
(\(Z=0\)), the problem reduces to estimating the relationship between the camera and this plane,
and then unprojecting the stabilized 2D trajectory onto the plane.

\medskip

To do so, we need the following information (all extracted on the reference frame):
(i) geometric cues that belong to the motion plane, such as two (or more) sets of \emph{parallel lines}
in the real world (e.g., track edges, fence/rail lines, painted markings), which allow estimating the
\emph{vanishing points} and therefore the \emph{vanishing line} \(l_\infty\) of the plane;
(ii) either the camera intrinsics \(K\) (from calibration) \emph{or} additional metric constraints
(orthogonality on the plane, known aspect ratio, or a known distance) to upgrade the rectification
from affine to metric.

\medskip

Given the vanishing line \(l_\infty\) and \(K\), the normal of the trajectory plane in camera coordinates
can be recovered up to scale (in particular, \(n \propto K^{-T} l_\infty\)).
This directly provides the \emph{angle between the camera and the plane} by comparing \(n\) with the
camera optical axis \(e_z\):
\[
\theta_{\text{plane}} = \arccos\!\left(\frac{n^\top e_z}{\|n\|}\right).
\]
After estimating the plane geometry, we rectify the reference frame with a plane-induced homography
\(H_{\text{rect}}\) that maps the plane to a fronto-parallel (metric) coordinate system.
In practice, this homography can be obtained either (a) from point correspondences between the image
and a planar layout with known geometry, or (b) by enforcing the vanishing-line constraint for affine
rectification and then applying metric constraints to obtain a Euclidean rectification.
The angle of the trajectory on the plane is given by the iformation on the slope.
\medskip

Finally, each original measurement \(p_t\) is mapped to plane (metric) coordinates by composing
stabilization and rectification:
\[
\tilde{p}_t \sim H_{\text{rect}}\,H^{\text{ref}}_t\,p_t,
\]
where \(H^{\text{ref}}_t\) maps frame \(t\) into the reference frame, and \(H_{\text{rect}}\) removes perspective
on the motion plane. The resulting trajectory is expressed in a planar metric coordinate system, and
the 3D trajectory is obtained by lifting it as \((X_t, Y_t, 0)\) under the \(Z=0\) assumption (with an
optional global scale set by one known distance on the plane).
